\subsection{Análisis de resultados y conclusiones}
\label{subsec:analisis-de-resultados-y-conclusiones}

\subsubsection{Práctica 1 - Densidad del Cilindro}

\begin{table}[H]
    \centering
    \caption{Densidades obtenidas por cada método y valores de referencia}
    \label{tab:comparacion-densidad}
    \begin{tabular}{|l|c|c|c|c|}
        \hline
        \textbf{Método} & $\rho$ ($\frac{g}{cm^3}$) & $\Delta\rho$ ($\frac{g}{cm^3}$) & \textbf{Lím. inf.} & \textbf{Lím. sup.} \\ \hline
        \csvreader[
            separator=semicolon,
            head=true,
            late after line=\\ \hline
        ]{secciones/analisis/comparacion-densidad.csv}{}
        {\csvcoli & \csvcolii & \csvcoliii & \csvcoliv & \csvcolv}
    \end{tabular}
    \begin{tablenotes}
        \small
        \centering
        \item[1] Densidades de referencia obtenidas de \cite{kayelaby1911}, pág. 20.
    \end{tablenotes}
\end{table}

\textbf{Compatibilidad entre métodos:} Los métodos son compatibles entre sí.
Los tres intervalos tienen intersección no nula, todos contienen el intervalo $[8{,}275; 8{,}405]$ $\frac{g}{cm^3}$.
La regla y la probeta tienen intervalos mucho más amplios debido a la menor precisión de esos instrumentos.

\textbf{Comparación con materiales de referencia:} La densidad obtenida con el calibre (método más preciso) es $\rho = 8{,}34 \pm 0{,}057 \frac{g}{cm^3}$.
La densidad de referencia más consistente es la de la constante de las pesas de latón $8{,}4 \frac{g}{cm^3}$.

Dado el aspecto del cilindro y los resultados revelados de su densidad, se estima que el mismo es una pesa patrón de laboratorio.
Si bien nuestra medición más precisa contiene la constante de la densidad utilizada para estos instrumentos ($8{,}4 \frac{g}{cm^3}$), lo hace en el límite superior del intervalo de indeterminación, como se puede apreciar en la Figura \ref{fig:grafico-intervalos}.

Se sabe que estos instrumentos están fabricados bajo especificaciones técnicas muy detalladas y precisas para garantizar las propiedades físicas.
Uno esperaría que nuestra medición más precisa se encuentre más próxima al centro del intervalo; estimamos que los errores introducidos por la medición no deberían ser lo suficientemente grandes para desplazar tanto el rango.

Nuestra hipótesis es que, como se muestra rudimentariamente en el esquema de la Figura \ref{fig:diagrama-cilindro}, el cilindro no es simplemente conexo, es decir, posee una cavidad interna.
Esto produce que nuestra fórmula de volumen para un cilindro perfecto resulte perjudicada, provocando que, al calcular la densidad ($\frac{M}{V}$), el denominador esté sobredimensionado y produzca el desplazamiento del rango antes mencionado.
Un mejor resultado se podría conseguir midiendo la altura y el diámetro de esta cavidad interna, calculando su volumen y restándole este al volumen total del cilindro.

\subsubsection{Práctica 1 -Análisis de Fuentes de Error}

\begin{itemize}
    \item \textbf{Error en la colocación de la regla:} Al no colocar la regla en contacto directo con el cilindro o al leer en ángulo, se introduce un sesgo sistemático.
    \item \textbf{Error de tension superficial en la probeta:} El líquido puede formar una curvatura provocada por la tension superficial con las paredes del instrumento, esto dificulta la lectura.
    \item \textbf{Burbujas de aire en la probeta:} Si quedó aire adherido al cilindro, el volumen desplazado es mayor al real, subestimando la densidad.
\end{itemize}

\subsubsection{Práctica 1 - Conclusiones}

Se determinó exitosamente la densidad del cilindro por tres métodos independientes, obteniendo resultados compatibles entre sí.
El calibre proporcionó la medición más precisa, mientras que la regla y la probeta tuvieron errores relativos mayores, evidenciando la importancia de elegir el instrumento adecuado según la precisión requerida.

% ####################################################################################################################
% ####################################################################################################################

\subsubsection{Práctica 2 - Pendulo Ideal}

\subsubsection{Resultado de $g$ y Comparación con Valor de Referencia:}

Tras realizar un ajuste con 12 puntos (correspondientes a $\theta$ = 5$\textdegree$, 9$\textdegree$ y 15$\textdegree$), y aplicando el error de la pendiente del ajuste, se obtuvo como resultado $g = 9{,}579 \pm 0{,}255 \frac{m}{s^2}$.

\textbf{Valor de referencia en Buenos Aires:} $g_{\text{ref}} = 9{,}797 \frac{m}{s^2}$.

El valor de referencia se encuentra dentro del intervalo de indeterminación del resultado obtenido ($9{,}834 \frac{m}{s^2}$, $9{,}324 \frac{m}{s^2}$) lo que indica que el resultado es compatible con el valor de referencia.

\subsubsection{Práctica 2 - Análisis de Fuentes de Error}

\paragraph{Práctica 2}
\begin{itemize}
    \item \textbf{Movimiento no plano (movimiento cónico):} si la esfera oscila en un cono en lugar de un plano, el período medido difiere del péndulo plano.
    \item \textbf{Masa no puntual:} El modelo asume una masa puntual, pero se utilizó una esfera de 1,7 $cm$ de radio.
    \item \textbf{Amplitudes grandes:} Para ángulos mayores a 5$\textdegree$ (especialmente el de 15$\textdegree$ incluido en el cálculo final), la aproximación $\sen \theta \approx \theta$ pierde validez, introduciendo un error sistemático.
\end{itemize}

\subsubsection{Práctica 2 - Conclusiones}

Se logró determinar de forma experimental la aceleración de la gravedad mediante el estudio de un péndulo simple.

A lo largo del desarrollo, se exploró no solo la obtención de un valor numérico para la aceleración de la gravedad, sino también la validez de los procesos aplicados.
Se propusieron y discutieron diversos métodos y criterios para el cálculo a partir de los datos, lo que nos permitió evaluar diferentes resultados junto con sus implicaciones.

Si bien el primer método de cálculo para la incerteza basado en mediciones individuales arrojó un error menor ($1{,}29\%$), se concluyó que el método basado en el error de la pendiente de la regresión lineal ($2{,}67\%$) es más apropiado para representar la calidad del ajuste al modelo teórico.

Para concluir, ambas prácticas ilustraron la importancia de la propagación correcta de incertezas, la elección del valor representativo, y la expresión de los resultados como $valor \pm incerteza [unidades]$.
